% !PRINTED none
\begin{center}
\gesperrt{Sechzehnter Abschnitt.}

\textbf{Zahlen und Functionale eines algebraischen K\"orpers.}

\rule{4cm}{0.4pt}
\end{center}

\begin{center}
\S.\,132.

\gesperrt{Definition der algebraischen Zahlen.}
\end{center}

Eine algebraische Gleichung
\[
F(x) = x^n + a_1 x^{n-1} + \cdots + a_{n-1} x + a_n = 0,
\]
deren Co\"efficienten $a_1, a_2, \ldots, a_n$ rationale Zahlen sind, nennen wir der K\"urze wegen eine rationale Gleichung. Sie hat, wie wir in fr\"uheren Abschnitten nachgewiesen haben, immer $n$ oder weniger, aber immer wenigstens eine Wurzel. Wie man jeder dieser Wurzeln durch rationale Zahlen, etwa durch Decimalbr\"uche oder durch Kettenbr\"uche, n\"othigenfalls mit Zuziehung der imagin\"aren Einheit $i = \sqrt{-1}$ bis auf jeden beliebigen Grad nahe kommen kann, d.\,h. wie man die Werthe der Wurzeln ann\"ahernd berechnen kann, ist im zweiten Buche des ersten Bandes gezeigt.

In den folgenden Betrachtungen soll es sich nun nicht um diese numerischen Werthe handeln, sondern um die arithmetischen Gesetze, denen diese Zahlen unterworfen sind, die sich aus der Definition selbst und nicht aus den numerischen Werthen ableiten lassen. Wir stellen also jetzt folgende Definition an die Spitze:

\gesperrt{Eine Zahl $\Theta$, die einer rationalen Gleichung}
\weq{1}{F(\Theta) = 0}
\gesperrt{gen\"ugt, heisst eine algebraische Zahl.}

Jede algebraische Gleichung mit rationalen Co\"efficienten liefert uns solche algebraische Zahlen, die sich also in beliebiger Menge angeben lassen. Die Frage, ob es auch nicht algebraische Zahlen giebt, wird uns sp\"ater besch\"aftigen.
